\section{Theorie}

\subsection{Das Periodensystem der Elemente}
Das moderne Periodensystem der Elemente (PSE) ordnet die chemischen Elemente in aufsteigender Reihenfolge basierend auf ihrer Ordnungszahl, welche der exakten Anzahl der Protonen im Atomkern und im elektrisch neutralen, elementaren Zustand auch der Anzahl der Elektronen in der Atomhülle entspricht. Der Aufbau des PSE spiegelt die Elektronenkonfiguration wider. Die Zeilen, werden als Perioden bezeichnet und fassen Elemente zusammen, deren Elektronen auf der gleichen Anzahl an Hauptschalen verteilt sind. Innerhalb einer Periode nimmt der Atomradius von links nach rechts ab, da die mit der Ordnungszahl zunehmende Kernladung zu einer stärkeren elektrostatischen Anziehung der Elektronen führt, wodurch die Elektronenhülle komprimiert wird. Die Spalten definieren die Elementgruppen, in denen die Elemente eine übereinstimmende Elektronenbesetzung in ihren äußersten \textit{s}- und \textit{p}-Orbitalen aufweisen. Dies führt zu ähnlichen chemischen und physikalischen Eigenschaften innerhalb einer Gruppe. Der Atomradius nimmt innerhalb einer Gruppe von oben nach unten zu, da mit jeder neuen Periode eine weitere, energetisch höhere Elektronenschale besetzt wird. \cite{skript}

\subsection{Die Orbitale}
Gemäß der Quantenmechanik lässt sich der Aufenthaltsort eines Elektrons nicht deterministisch festlegen. Stattdessen beschreiben Orbitale als Lösungsfunktionen der Schrödinger-Gleichung dreidimensionale Räume, in denen die Aufenthaltswahrscheinlichkeit des Elektrons am größten ist. Form, Größe und Ausrichtung dieser Orbitale werden dabei formal durch drei Quantenzahlen definiert. Die Hauptquantenzahl $n$ bestimmt das Energieniveau, die Nebenquantenzahl $l$ die Geometrie der Orbitale und die Magnetquantenzahl $m_l$ die räumliche Orientierung. Die Besetzung dieser Orbitale erfordert zusätzlich die Spinquantenzahl $m_s$ des Elektrons und unterliegt drei zentralen Prinzipien. Dem Aufbauprinzip, welches die Besetzung ausgehend vom niedrigsten Energiezustand beschreibt. Der Hundschen Regel, die die maximale Anzahl ungepaarter Elektronen in entarteten Orbitalen festlegt, sowie dem Pauli-Prinzip. Letzteres fordert, dass keine zwei Elektronen in allen vier Quantenzahlen übereinstimmen dürfen. Befinden sich also zwei Elektronen im selben Orbital, müssen sie sich zwingend in ihrem Spin unterscheiden. \cite{skript}

\subsubsection{\textit{s}-Orbitale}
Das \textit{s}-Orbital stellt den energetisch tiefsten Zustand innerhalb einer Elektronenschale dar und weist eine kugelsymmetrische Form um den Atomkern auf. Ab der ersten Periode verfügt jede Hauptschale über genau ein \textit{s}-Orbital, das maximal zwei Elektronen fassen kann und keine Knotenflächen besitzt. \cite{skript}

\subsubsection{\textit{p}-Orbitale}
Ab der zweiten Periode treten \textit{p}-Orbitale auf, die eine hantel- oder keulenförmige Geometrie besitzen und sich räumlich entlang der drei Achsen orientieren. Es existieren drei energetisch gleichwertige \textit{p}-Orbitale pro Schale, die durch eine Knotenfläche im Atomkern getrennt sind und zusammen bis zu sechs Elektronen aufnehmen können. \cite{skript}

\subsubsection{\textit{d}-Orbitale}
Die \textit{d}-Orbitale sind ab der dritten Periode besetzbar, weisen deutlich komplexere räumliche Strukturen auf und bilden das höchste Energieniveau der hier betrachteten Orbitale. Von den insgesamt fünf \textit{d}-Orbitalen besitzen vier eine gekreuzte Doppelhantelform (kleeblattartig), während das fünfte hantelförmig mit einem äquatorialen Ring (Torus) strukturiert ist. Sie besitzen jeweils zwei Knotenflächen und können bis zu zehn Elektronen beherbergen. \cite{skript}

\subsection{Das HSAB-Prinzip}
Das HSAB-Prinzip welches für Hard and Soft Acids and Bases-Prinzip steht, bietet ein qualitatives Modell zur Vorhersage der Stabilität von Lewis-Säure-Base-Komplexen. Es kategorisiert Säuren und Basen anhand ihrer Härte beziehungsweise Weichheit. Diese hängt primär von ihrer Größe, Ladungsdichte und Polarisierbarkeit ab. Als hart gelten Teilchen mit einem geringen Volumen, einer hohen Ladung, einer hohen Elektronegativität und einer folglich niedrigen Polarisierbarkeit. Im Gegensatz dazu sind als weich bezeichnete Teilchen durch einen großen Radius, eine geringere Ladung, eine hohe Polarisierbarkeit und oft durch besetzte \textit{d}-Orbitale charakterisiert. \cite{skript}

Die fundamentale Kernaussage des Konzepts lautet, dass harte Lewis-Säuren bevorzugt stabile Bindungen mit harten Lewis-Basen eingehen, während weiche Lewis-Säuren stärkere Wechselwirkungen mit weichen Lewis-Basen präferieren. Kombinationen aus harten und weichen Reaktionspartnern führen aufgrund der energetischen und sterischen Diskrepanzen zu deutlich instabileren Verbindungen. \cite{skript}

\newpage